\documentclass[11pt]{article}
\usepackage[margin=1in]{geometry}
\usepackage{amsmath,amssymb,amsthm}
\usepackage[colorlinks=true,urlcolor=blue,linkcolor=blue]{hyperref}
\newtheorem{theorem}{Theorem}
\newtheorem{lemma}[theorem]{Lemma}
\newcommand{\cost}{\operatorname{cost}}
\newcommand{\OPT}{\operatorname{OPT}}
\newcommand{\MS}{\mathrm{MS}}
\newcommand{\polylog}{\operatorname{polylog}}
\setlength{\emergencystretch}{4em}
\title{A common forest interface for the splay comparison proof}
\author{Research note; hand extension of existing interfaces}
\date{2 October 2026}
\begin{document}
\maketitle

\noindent\textbf{Status.}
This note records a positive extension suggested by the existing Lean
normal-form interfaces. The forest composition is a hand argument, not
an additional machine-checked theorem. The final splay comparison is
conditional on the accounting argument of the current splay manuscript.
No bound on the cost of an appropriate multi-splay comparator is assumed
or proved here.

\section{The interface and its cost model}

At every time there is a forest of disjoint ordered key sets. Each set
has one BST and a root handle. An operation may select a component at
unit cost; this convention is shared by the splay algorithm and its
comparator. It does not grant a free move to an interior node. The
operation word specifies common component domains for both algorithms.
The allowed operations are:
\begin{enumerate}
\item \emph{Search}$(x)$: serve $x$ in its specified current component.
The splay algorithm splays $x$ to that component's root.
\item \emph{Split}$(x)$: serve and move $x$ to the root, then remove
$x$, returning the two ordered child domains as separate components.
The splay algorithm splays $x$ before this root deletion.
\item \emph{Join}$(L,x,R)$: take two components whose keys are all
below and above the supplied inactive pivot $x$, and replace them by
the tree $x(L,R)$. The comparator performs the same common root join;
it may subsequently rotate its tree. The splay algorithm initially
uses exactly the displayed root join, without rebuilding the children.
\end{enumerate}
Empty children are permitted. A removed pivot can be supplied again
later. Let $N$ be the number of initially active keys, $q$ the number of
structural operations, and $K$ an upper bound on every live component
size, including immediately after a join. Let $C$ be a comparator's
total search, pointer, and rotation cost; charging the structural
operations separately as $q$ avoids any ambiguity about their unit
cost convention.

The main distinction from independently restarting an analysis is that
the actual representation and accounting state of each surviving child
travel with that child. A join combines the existing states; it does
not initialize a new analysis on all keys in the new component.

\section{Existing local normal-form interfaces}

Write $L_K=1+\lceil\log_2(K+1)\rceil$. The repository has two
representation layers with exactly the required root operations.

\paragraph{The heavy-path layer.}
\texttt{OfflineHeavyJoin.lean} proves
\texttt{HeavyRepresentation.exists\_join};
\texttt{OfflineHeavySplit.lean} proves
\texttt{exists\_split\_node}. For either operation there is an exposed
supported path of cost $c_H$ satisfying
\[
 c_H\le 16L_K^2+3,
 \qquad
 c_H+\Phi_H(\text{after})
 \le\Phi_H(\text{before})+115.
\]
Before and after potentials are summed over the participating children
when the state is a forest. Every represented tree has height at most
$6\lceil\log_2(K+1)\rceil^2$. The potential bound
\texttt{OfflineHeavyPotential.lean} is
\[
 0\le\Phi_H(T)\le64|T|.
\]
The split theorem expressly returns the child representation histories.

\paragraph{The rib layer.}
\texttt{OfflineRibRootOperations.lean} proves a root join and root
split with
\[
 c_R\le64L_K^2,
 \qquad
 c_R+\Phi_R(\text{after})
 \le\Phi_R(\text{before})+448.
\]
A join requests a left representation with orientation \texttt{true}
and a right representation with orientation \texttt{false}.
The same file proves \texttt{exists\_reorient} for either desired
orientation, with
\[
 c_{\mathrm{orient}}\le128L_K^2,
 \qquad
 c_{\mathrm{orient}}+\Phi_R(\text{after})
 \le\Phi_R(\text{before})+896.
\]
This operation preserves the represented source tree. It consists of
two root stack operations, rather than a reconstruction of the component.
\texttt{OfflineRibPotential.lean} gives
$0\le\Phi_R(T)\le256|T|$.

Rib representations have bounded minor ribs, with slack nineteen, and
increase source height by at most a factor
$5\lceil\log_2(K+1)\rceil$.
The supported-path conversion in
\texttt{OfflineRestrictedMacros.lean} and
\texttt{OfflineRibbedTrace.lean} replaces each supported macro by
root or grandchild rotations at constant overhead in its support size.
It preserves a height bound equal to the endpoint height plus the
support size, and a constant minor-rib bound (slack twenty for the
ribbed construction).

The underlying representation, rotation, split and join statements are
parameterized by a \emph{capacity}. It is therefore legitimate to use the
same fixed capacity $K$ throughout an evolving forest. Changing a
component's size does not require changing its representation parameter.
The exported \texttt{exists\_boundedService} and
\texttt{exists\_boundedRotation} corollaries specialize that capacity to
the dense alphabet size. Here their proofs are applied with arbitrary
capacity $K$: the displayed support and potential inequalities come from
the underlying capacity-parameterized interfaces, and their conversion
to shallow traces is unchanged. This lifting is part of the present
hand argument, not an existing forest theorem in Lean.

A live component's possibly nonconsecutive keys can be relabeled by
their inorder positions for these local statements. Splitting or joining
changes this order isomorphism but changes no physical tree shape,
rotation, stack digit, supported frontier, or potential. Transport the
existing representation witnesses through the new order isomorphism;
do not choose fresh representations. Thus alphabet relabeling itself
performs no pointer operations and grants no access to an interior node.

\section{The forest normal form}

\begin{theorem}[Common forest normalization]\label{thm:normal}
Every comparator execution of the interface above has a common-domain
comparison forest with total cost
\[
 O(C+N+q),
\]
height $O(L_K^3)$ in every physical intermediate tree, a constant bound
on minor ribs, and all rotations at constant depth within the component
being processed. Immediately before a split and immediately after the
common raw join, the relevant pivot is the root. All surviving child
representation histories are retained across the structural operations.
\end{theorem}

\begin{proof}
First replace each root-starting comparator transaction by shallow
rotations of its touched rooted hull, retaining the same hanging
subtrees. This is the usual constant-cost restricted-rotation conversion,
applied inside the selected component. A split first exposes its
requested pivot as a root; a join is a raw common root insertion followed
by any further comparator work. This gives $O(C+q)$ root services,
shallow rotations, root deletions and root insertions.

Maintain a heavy representation of each component at fixed capacity
$K$. The supported rotation and service interfaces already used in
\texttt{OfflineHeavyBoundedSteps.lean} handle the fixed-domain work.
Use the explicit heavy split and join interfaces for the structural
operations. The raw join $x(A,B)$ has height at most one more than the
maximum child height, so it also satisfies an $O(L_K^2)$ bound.
The supported paths have size $O(L_K^2)$; conversion to shallow rotations
therefore keeps all intermediate heights $O(L_K^2)$.

Sum the heavy potential over all live components. A split replaces one
summand by its two surviving child summands, and a join performs the
reverse accounting operation with a new pivot. The initial potential is
at most $64N$. Each primitive supplies only its stated constant
amortized allowance. Nonnegativity of the final global potential gives
$O(C+q+N)$ total work in this first representation layer. There is no
component-lifetime initialization sum.

Apply the rib representation to the resulting shallow trace, again
retaining its histories at fixed capacity $K$. Before a raw root join,
reorient the two child representations as necessary. There are at most
two such reorientations. Insert the common pivot above the children,
then apply the rib join interface. Its source has bounded minor ribs
because of the chosen child orientations; this is the same exposed-root
geometry used by \texttt{exists\_boundedService} in
\texttt{OfflineRibBoundedSteps.lean}. Its height is one more than the
maximum child height. A root split is handled by exposing the source
root with \texttt{exists\_split\_node}, deleting that root, and retaining
the two returned child representations. Their returned orientations may
be kept until a later join requires a change.

The rib height is $O(L_K)$ times the first layer's $O(L_K^2)$ height;
the $O(L_K^2)$ supported-path allowance is smaller. Supported-path
conversion keeps this $O(L_K^3)$ height and constant minor-rib bound
throughout every macro. Its rotations stay at constant depth.

Finally sum the rib potential over live components. Its initial sum is
at most $256N$. A root join costs at most two reorientation allowances
and one join allowance; a split costs one split allowance. The existing
rotation and service allowances charge the rest of the first-layer
trace. The potentials telescope with the actual child histories.
Thus the second layer costs $O(C+q+N)$ as well.

For synchronization, expose the comparison pivot before the algorithm's
splay for a split, and delete the common root only afterwards. At a
join, perform any orientation changes while the children are still
separate; then both forests receive the same raw root $x$. Subsequent
comparison normalization rotations are ordinary paid comparator changes.
This establishes the stated common-root checkpoints.
\end{proof}

This proof is an application of the local existing interfaces. A Lean
formalization would additionally need a forest trace type, the global
sum of component potentials, and order-isomorphism transport between
component alphabets. These are not supplied by the existing fixed-domain
\texttt{RibbedTrace} datatype.

\section{Root joins in the splay accounting state}

\begin{lemma}[Common root join accounting]\label{lem:heap-join}
Suppose the splay and comparison forests join the same two domains
under a new common root $x$. The lazy-interval accounting state of the
splay manuscript can retain all old states and perform only $O(1)$ new
interface and potential events for this operation.
\end{lemma}

\begin{proof}
Immediately before the join, ranks in each child are measured from
the root of its own comparison component. After the common raw join,
increase all old ranks in each child by one and give $x$ rank zero.
All differences between old ranks within a child are unchanged.
Consequently all internal gaps, point gaps, heavy paths and lazy
intervals of that child retain their previous state.

The new minimum $x$ becomes the root heavy path. The previous root
heavy path of each nonempty child becomes one left or right heap child
of $x$. Its minimum comparison rank used to be zero and is now one,
so its new gap is one. Create the two possible growing singleton
intervals at $x$ and their constant-size records. This introduces only
$O(1)$ point-gap contributions and boundary events, and only $O(1)$
new bends. All records internal to the children remain in place.

More explicitly, the old child-root has interface value $f=1$ from
its root point gap zero. Its new singleton gap one has contracted
value $f=2$. Apply one \texttt{ValueIncrease} to this child root,
\texttt{Create} the corresponding growing parent slot at value one,
and \texttt{Insert} the child at value two. The insertion inequality
is $2\ge1+1$, and there are at most two such children. The inverse
root deletion first \texttt{Convert}s these root singletons, returning
the child's value from two to one, and then deletes them. Empty
growing slots of the removed pivot are converted and retired. On
reinsertion its newly created slots again start at value one.
Child intervals, pass counts, and creator histories are retained.

This is the reverse of the common root deletion in
\texttt{app-split-conjecture.tex}, except that insertion is charged its
constant number of new records. Reinserting a formerly deleted pivot
creates those records again and is paid again as one of the $q$
structural operations. It does not recreate the children's records.
One can place every key that will occur into the interface universe
in advance, giving an inactive key the neutral value $f=1$ and
charging the initial excess $\sum(f-1)$. Alternatively, its size is
at most the $N$ initially active keys plus the $q$ insertion events.
Either convention charges future pivot lifetimes by $O(N+q)$ and
does not reset a live child's potential.
Subsequent comparison normalization rotations are handled by the
ordinary paid-rotation part of the original proof.
\end{proof}

\begin{theorem}[Conditional common-forest splay comparison]
Assuming the core lazy-interval and pairing-accounting argument of the
current splay manuscript, its comparison extends to the interface above
with
\[
 \cost_{\mathrm{splay\ forest}}
 =O\!\left((C+N+q)\,
   (1+\log\log(K+4))
   (1+\log\log\log(K+16))^2\right).
\]
\end{theorem}

\begin{proof}
Use Theorem~\ref{thm:normal} as the comparison forest. Its ranks lie in
$O(L_K^3)$, so the maximum contracted rank gap is $O(1+\log\log(K+4))$.
The two logarithmic losses from the existing lazy-interval accounting
are $O(1+\log\log\log(K+16))$ each.

The original forest has $O(N)$ initial records. Ordinary splays and
paid comparison rotations create the already established constant
number of records per event. Common root deletion uses the unchanged
split argument; Lemma~\ref{lem:heap-join} handles common root insertion.
Thus all record and boundary budgets are $O(C+N+q)$, and the original
global pairing and bend estimates sum over the live components.
No potential is restarted at a split or join.

Here is the global accounting explicitly. Set $B_0=N+C+q$ and let
$G\ge2$ bound the contracted gaps. Let $\ell$ count created lazy
intervals, $p$ paid interface operations, $u$ unimportant boundary
pairings, $g$ good internal pairings, and $A$ all internal pairings.
The retained-history construction and the local splay estimates give
\[
 \ell=O(B_0),\qquad
 p\le u+O(B_0\log G),\qquad
 g\le uG+O(B_0G\log G).
\]
Root joins add only $O(q)$ inserts and value increases. Thus the
interface totals satisfy
\[
 \#\texttt{Insert}=O(p+N+q),\qquad
 \#\texttt{ValueChange}=O\bigl(g+(p+N+q)G\bigr).
\]
The original pairing inequality consequently remains
$A\le\alpha(pG\log G+O(B_0G\log^2G))+O(1)$ for its fixed constant
$\alpha$. The same definition of unimportance supplies
$\alpha uG\log G\le A/2$. Absorbing this term gives
$A=O(B_0G\log^2G)$. Important boundary pairings contribute
$O(B_0\log G)$, and unimportant ones are absorbed as before. Joins
and splits add $O(q)$ bends, so the zig-zag count has the same bound.
Taking $G=O(1+\log\log(K+4))$ gives the displayed theorem.
\end{proof}

\section{Optional historical compression}

The same forest construction can be fed into the later strict-interface
compression argument in
\texttt{research/approximation-2026-09-11/}.
This is an optional hand corollary of the forest realization above;
it is not a separately formalized forest theorem. In the notation of
that argument, let $B$ now denote all boundary pairings and retain
$B_0=N+C+q$. The unchanged global estimates are
\[
 I=O(B_0+B),\qquad V=O(G(B_0+B)),\qquad
 B\le O(B_0(1+\log\kappa))+O(A/\kappa)
\]
for every boundary cutoff $\kappa$. Root joins contribute only $O(q)$
inserts and value changes; neither their retained child records nor
root splits reset source pass counters. Thus the historical-rank
compression theorem applies to this same single strict trace.

Write $L_j(t)$ for stabilized $j$-fold iterated logarithms. One
compression, as in \texttt{HistoricalCompressionIteratedLog.lean},
replaces the displayed forest factor by
\[
 g_1(K)=O\bigl(L_2(K)L_3(K)L_4(K)^2\bigr).
\]
Two applications of the literal coefficient transfer in
\texttt{HistoricalCompressionRecursion.lean} and its concrete boundary
absorption give the fixed-depth hand asymptotic simplification
\[
 g_2(K)=O\bigl(L_2(K)L_3(K)L_4(K)L_5(K)^2\bigr).
\]
The recursion's constants must remain attached to its chosen depth;
these statements use a fixed depth of one or two, not a depth growing
with $K$. For example, combining the two-compression version with a
forest comparator of factor $\lambda=O(L_2(H))$, where
$H=\Theta(\log n)$, yields
\[
 O\bigl((\OPT(X)+n)L_3(n)^2L_4(n)L_5(n)L_6(n)^2\bigr).
\]
The leading squared $L_3(n)$ remains. This refinement therefore does
not resolve the exact $O(L_3(n))$ multi-splay question.

\section{An exact common forest word for multi-splay}

Consider a left-to-right primary switch at $y$. Its current splay
component contains $U\cup\{y\}\cup L$, where $U$ is the preferred
prefix above $y$ and $L$ is the old preferred tail in the left child of
$y$. The incoming right tail $R$ is a separate splay component. Let
$\ell<y<r$ be the closest lower and upper ancestors in $U$, when
present. These are the two guards used by standard multi-splay.

Assume first that both guards are present. The exact operation word is
\[
 \operatorname{Split}(y),\quad
 \operatorname{Split}(\ell),\quad
 \operatorname{Split}(r),\quad
 \operatorname{Join}(A,\ell,\varnothing),\quad
 \operatorname{Join}(R,r,B),\quad
 \operatorname{Join}(A',y,B').
\]
The second split selects the left daughter of the first split, and the
third selects its right daughter. Here $A$ is the part of $U$ below
$\ell$, $B$ is the part above $r$, and $A',B'$ are the outputs of the
two preceding joins. The right daughter of the $\ell$ split is $L$;
it remains parked. The left daughter of the $r$ split is empty because
the current owned tail lies to the left of $y$.

Splitting at $y$ performs its full splay. The guard splits perform
exactly the splays stopped at the corresponding child of $y$; removing
$y$ analytically makes those stopped splays ordinary root splays of
the daughter components. The three raw joins restore the guards and
$y$ in exactly the shape produced by the standard switch. Marking the
old tail and unmarking the incoming tail are precisely the corresponding
changes of component membership. Foreign components in the inorder
gaps remain opaque throughout.

If $\ell$ is absent, the first split's entire left daughter is the
parked tail $L$ and the new left daughter is empty. If $r$ is absent,
the new right daughter is simply $R$. Omit the absent guard's split and
join. Right-to-left switches use the reflected operation word. The
same construction applies to the final switch at the requested key;
missing children merely give empty tails.

Consequently there is one fixed mixed forest word $W_{\MS}(X)$,
determined by the reference tree, preferred directions and original
request sequence, with at most six structural operations per switch.
It is independent of the shapes of the auxiliary splay trees. Every
intermediate component has at most $H$ keys, where $H$ is the maximum
preferred-path length. The actual forest splay execution of this word
has exactly the three possible switch splays and their resulting
auxiliary shapes. The original multi-splay search traversal is charged
to these splays by the standard traversal argument; thus its physical
cost is at most a constant times their cost plus $|X|+S_{\MS}$.

For completeness, this traversal charge is pointwise. In each initial
auxiliary exited by the original search, both neighboring keys of the
exit gap lie on its search path. One is the subsequent switch pivot
and the other is one of its ancestor guards. Bottom-up work in deeper
components changes only the incoming opaque payload, so this auxiliary
has not changed before its switch. If the pivot is the deeper of the
two gap neighbors, its first splay already pays the entire old gap
search. Otherwise the first splay pays the path to the pivot, and the
stopped guard splay pays the surviving path below it to the other
neighbor. The final initial auxiliary contains the requested key.
Ancestor switches attach it as an opaque payload, preserving its
internal root-to-query path until the mandatory final query splay.
That last splay pays the remaining initial search segment. Constant
cost per auxiliary boundary is included in $|X|+S_{\MS}$.

Let $C_{\mathrm{forest}}(W_{\MS}(X))$ be the cost of a comparator for
this same complete mixed word. The conditional common-forest theorem
therefore gives the concrete reduction
\[
 \cost_{\MS}(X)
 =O\!\left((C_{\mathrm{forest}}(W_{\MS}(X))+n+S_{\MS})\,g(H)\right),
\]
where
\[
 g(H)=(1+\log\log(H+4))
      (1+\log\log\log(H+16))^2.
\]
The inherited child histories in the normal-form proof are essential
for this reduction: a new analysis is not started after each of the
six forest operations.

\section{The remaining comparator transport}

Using $K=O(\log n)$ would give the factor
\[
 O(\log\log\log n\;(\log\log\log\log n)^2).
\]
The exact word above leaves the following substantive task: exhibit a
comparator for that same word of cost
$O(\OPT(X)+n)$, or a suitably explicit weaker bound.

Neither an independently optimal tree for each path lifetime nor the
query-only cost of a secondary Tango instance supplies this comparator.
The forest word prescribes its splits and joins as well as its searches.
Conversely, the normal-form and initialization problems for the precise
common root split/join interface need not be treated as an unexplained
additional conjecture: the existing representation inequalities above
give a concrete route to their solution.

For example, a compatible comparator with cost
$O((\OPT(X)+n)(1+\log\log H))$ would yield, through this reduction,
the weaker factor
\[
 O\!\left((\log\log\log n)^2
          (\log\log\log\log n)^2\right).
\]
Such a comparator must support the prescribed guard exposures and
root interfaces; the ordinary query-only Tango-squared bound does not
by itself certify that property. This note does not assert the
existence of either comparator bound.

\end{document}
